\chapter*{Notes}
\label{ch:notes}
\markboth{Notes}{Notes}
\addcontentsline{toc}{chapter}{Notes}

\section*{Notes on Chapter I}

\noindent For an introduction to Fourier analysis on groups, see the classic by Weil~\cite{weil1940}; also Loomis~\cite{loomis1953}, Naimark~\cite{naimark1964}, the Appendix in Goldberg~\cite{goldberg1961}, Rudin~\cite{rudin1962}, and Reiter~\cite{reiter1968}, where the contributions of A.~Beurling, I.E.~Segal, and others, are described.

\noindent\textbf{§3.}\quad For the theory of distributions, in general, with applications, see the classics by L.~Schwartz~\cite{schwartz1950}, and I.M.~Gelfand and G.E.~Shilov~\cite{gelfandshilov1964}. For distributions in connexion with Fourier transforms, in particular, see Ch.\,I of H{\"o}rmander's book~\cite{hormander1963}, also Yosida~\cite{yosida1965}, Ch.\,VI, and Donoghue~\cite{donoghue1969}.

\noindent The function~$\omega$ in~\eqref{eq:I-3.12} was introduced by Wiener~\cite{wiener1926}, p.\,562.

\noindent\textbf{§4.}\quad The localization theorem here is motivated by the one on Fourier series due to Riemann, see Hardy and Rogosinski~\cite{hardyrogo1944}, pp.\,39--42, and Zygmund~\cite{zygmund1959}, Ch.\,II, §6, p.\,53, §8.

\noindent The examples of Mellin transforms given here are of frequent occurrence in analytic number theory, see, for instance, the author's book~\cite{chandrasekharan1970}. Several more are given by Titchmarsh~\cite{titchmarsh1937}.

\noindent\textbf{§5.}\quad For Poisson's summation formula, see Bochner~\cite{bochner1932}; his proof is also given in the author's book~\cite{chandrasekharan1970}. For some special applications see, for instance, Zygmund~\cite{zygmund1959}, Vol.\,I, Ch.\,II, §13.

\noindent For the theta-relation~\eqref{eq:I-5.8}, in the general setting of theta-functions, see, for instance, the author's book~\cite{chandrasekharan1985}, where the connexion with elliptic functions, and the theory of numbers, is elucidated on an elementary level.

\noindent There is also an~$L_{2}$-version in one variable, see Boas~\cite{boas1946}.

\noindent\textbf{§6.}\quad The proof of the uniqueness in Theorem~\ref{thm:I-7} (without the use of summability and general inversion) can be effected by the use of a piece-wise linear (trapezoidal) function instead of the function~$\omega_{c,\varepsilon}$; see Bochner and Chandrasekharan~\cite{bochnerchandrasekharan1949}, Ch.\,I, §6, Th.\,5.

\noindent For a sharper version of the uniqueness theorem due to A.C.~Offord~\cite{offord1936}, see~\ref{thm:III-10}.

\noindent\textbf{§7.}\quad The motivation for the summability theorems here is again supplied by the theory of Fourier series, see, for instance, Hardy and Rogosinski~\cite{hardyrogo1944}, Ch.\,V, p.\,70; Bochner and Chandrasekharan~\cite{bochnerchandrasekharan1949}, §7.

\noindent For properties~\eqref{eq:I-7.6} and~\eqref{eq:I-7.20} of integrable functions, and for the definition of the ``Lebesgue set'', see, for instance, Titchmarsh~\cite{titchmarsh1939}, §11.6.

\noindent Convolution integrals of the type~\eqref{eq:I-7.1} are of importance in the theory of approximation. See Butzer and Nessel~\cite{butzernessel1971}, where generalized singular integrals of the type of Cauchy-Poisson, Gauss-Weierstrass, Fej\'er, and Bochner-Riesz, are dealt with in detail. For more general methods, see Stein~\cite{stein1970}, Stein and Weiss~\cite{steinweiss1971}.

\noindent\textbf{§8.}\quad Example~3, following Theorem~\ref{thm:I-11}, is used by C.L.~Siegel in his proof of Hamburger's theorem on the Riemann zeta-function. See, for instance, the author's book~\cite{chandrasekharan1970}, Ch.\,II, §5.

\noindent Theorem~\ref{thm:I-12} is due to Bochner and Chandrasekharan~\cite{bochnerchandrasekharan1949}, Ch.\,I, Th.\,9, p.\,20; also p.\,211, where it is commented upon. A generalization was later given by Bochner in his book~\cite{bochner1955}, p.\,25, Th.\,2.2.1.

\noindent Theorems~\ref{thm:I-13} and~\ref{thm:I-14} make it possible to define the Fourier transform on~$L_2(-\infty,\infty)$ and prove Plancherel's theorem [cf.~Ch.\,II] by starting from the subspace~$L_1\cap L_2$.

\noindent\textbf{§9.}\quad The systematic use of summability in norm seems to have originated with Wiener. For Theorem~\ref{thm:I-17} see Bochner and Chandrasekharan~\cite{bochnerchandrasekharan1949}, Ch.\,I, §10, who also proved further results in that direction.

\noindent On Weyl's form of the Riesz-Fischer theorem, see Weyl~\cite{weyl1909}, and Wiener's~\cite{wiener1933} remarks; also Stone's~\cite{stone1932}, p.\,26; and J.~von Neumann's~\cite{neumann1929}, p.\,109--111.

\noindent\textbf{§10.}\quad For the central limit theorem, see Cram\'er~\cite{cramer1946}, Ch.\,17, §4, and Feller~\cite{feller1950}, Ch.\,VIII, §4. According to Cram\'er, the theorem was first stated by Laplace in 1812; a rigorous proof under ``fairly general'' conditions was given by Liapounoff; and the problem of finding the most general conditions of validity was solved by Feller, Khintchine, and L\'evy~\cite{levy1934}. The proof given here differs only in detail from that given, for instance, by Dym and McKean~\cite{dymmckean1972}, Ch.\,2, §7.

\noindent\textbf{§11.}\quad Theorem~\ref{thm:I-21} is the analogue, for Fourier transforms, of a classical theorem on the absolute convergence of Fourier series due to N.~Wiener~\cite{wiener1930} and P.~L\'evy~\cite{levy1934}. See Zygmund~\cite{zygmund1959}, Vol.\,I, Ch.\,VI, §5. The proof given here differs only in detail from that of R.R.~Goldberg~\cite{goldberg1961}, Ch.\,2, §9, which is itself closely modelled on Bochner's proof~\cite{bochner1936a} of the Wiener-L\'evy theorem.

\noindent\textbf{§12.}\quad Wiener was the first to study ``closure'' properties of functions in~$L_1(-\infty,\infty)$ and in~$L_2(-\infty,\infty)$, and relate them to Fourier transform theory. See Wiener~\cite{wiener1930}. The proof given here of Theorem~\ref{thm:I-23} is the same as Bochner's~\cite{bochner1936a}. An algebraic reformulation of the theorem would be that every proper closed ideal of~$L_1(-\infty,\infty)$ is contained in a maximal ideal. The problem of characterizing the sub-class of functions~$f$ in~$L_1(-\infty,\infty)$ which have the property that~$\overline{S}_f$ is the intersection of the maximal ideals containing it has received attention. For the work of Beurling and others, see Pollard~\cite{pollard1953}, and Reiter~\cite{reiter1968}. For generalizations of Theorem~\ref{thm:I-23}, see Ch.\,4 of Goldberg~\cite{goldberg1961}, Reiter~\cite{reiter1968}, where further references can be found, e.g. to Agmon and Mandelbrojt~\cite{agmonmandelbrojt1950}, Malliavin~\cite{malliavin1959}, and others.

\noindent\textbf{§13.}\quad Theorems~\ref{thm:I-24} and~\ref{thm:I-25} are due to Wiener~\cite{wiener1930}. Theorem~\ref{thm:I-26} is due to Littlewood~\cite{littlewood1910}, and forms the prototype for many of the tauberian theorems of the period before Wiener. Littlewood's theorem can be proved directly, and simply, as Karamata~\cite{karamata1930} has shown, by the use of Weierstrass's theorem on the approximation of continuous functions, instead of Wiener's theorem on the~$L_1$-closure. See Wiener's own remarks~\cite{wiener1933}, and Wielandt's~\cite{wielandt1952} arrangement of Karamata's proof.

\noindent Wiener's work on tauberian theorems has been carried forward notably by H.R.~Pitt~\cite{pitt1958}.

\noindent Albert Stadler~\cite{stadler1986} has recently proved a tauberian theorem, with remainder, of the Wiener-Ikehara type (see, for instance, the author's book~\cite{chandrasekharan1968}), which yields the more refined forms of the prime number theorem as corollaries. See Wiener's remarks~\cite{wiener1930}, p.\,93, on this possibility.

\noindent Bochner and Chandrasekharan~\cite{bochnerchandrasekharan1949}, Th.\,29, p.\,54, subsume Karamata's theorem as part of another theorem which characterizes what they call the Karamata extension of the kernel~$e^{-\alpha}$. The nature of this extension in the case of general kernels seems not to be known.

\section*{Notes on Chapter II}

\noindent\textbf{§14.}\quad Theorems~27 and~28 are due to Bochner and Chandrasekharan~\cite{bochnerchandrasekharan1949}, Ch.\,I, §15. See Butzer and Nessel~\cite{butzernessel1971} Ch.\,7, for later developments. Equations~\eqref{eq:I-14.1} and~\eqref{eq:I-14.2} arise in connexion with the problem of conduction of heat, see Carslaw~\cite{carslaw1921}, §§16,45.

\noindent\textbf{§15.}\quad The proof of~\eqref{eq:I-15.24} given here is the same as the one given by Bochner in his book~\cite{bochner1932}; he comments that according to Burkhardt~\cite{burkhardt1904} pp.\,1165--1173, the cases $k=2,3$ are due to Poisson and Cauchy, and that the ``theorem is also not new for $k$ arbitrary''. A second proof is given by Bochner and Chandrasekharan~\cite{bochnerchandrasekharan1949}, pp.\,71--74.

\noindent The introduction of the spherical mean~$f_{x}(t)$ is due to Bochner. He carried the idea further into the study of multiple Fourier series. See Bochner~\cite{bochner1936b}, followed by Chandrasekharan~\cite{chandrasekharan1948}, Chandrasekharan and Minakshisundaram~\cite{chandrasekharanminakshisundaram1947}, and~\cite{chandrasekharanminakshisundaram1952}, Ch.\,IV, and H.~Joris~\cite{joris1968}. Important work with quite different techniques has been accomplished on topics in multiple Fourier series by E.M.~Stein, and others. See, for instance, Ch.\,VII of Stein and Weiss~\cite{steinweiss1971}.

\noindent The evaluation of~$V_{k}(s)$ by induction is done, for instance, by Walfisz~\cite{walfisz1957}, p.\,41.

\noindent\textbf{§2, §3.}\quad Different proofs of Plancherel's~\cite{plancherel1910} theorem have been given by Titchmarsh~\cite{titchmarsh1939}, Bochner~\cite{bochner1932}, F.~Riesz~\cite{riesz1927}, Wiener~\cite{wiener1933}, M.H.~Stone~\cite{stone1932}, p.\,104, and Bochner and Chandrasekharan~\cite{bochnerchandrasekharan1949}. For the Remarks following~\eqref{eq:II-2.29} see Stein and Weiss~\cite{steinweiss1971}, p.\,18. For~\eqref{eq:II-3.5} see Bochner and Chandrasekharan~\cite{bochnerchandrasekharan1949}, p.\,99.

\noindent\textbf{§4.}\quad Theorem~\ref{thm:II-6} is due to Wiener; for Theorem~\ref{thm:II-7} see Bochner and Chandrasekharan~\cite{bochnerchandrasekharan1949}, Ch.\,IV, §10.

\noindent\textbf{§5.}\quad Weyl's proof of the inequality under somewhat stronger hypotheses is given in Appendix~I to his book~\cite{weyl1928}. His proof in the second edition differs in detail from the one given in the first. The proof given here differs only in detail from the one given in their book by Dym and McKean~\cite{dymmckean1972}.

\noindent\textbf{§6.}\quad The Phragm\'en-Lindel\"of~\cite{phragmenlindelof1908} principle takes many forms. See, for instance, Calder\'on-Zygmund~\cite{calderonzygmund1950}, Littlewood~\cite{littlewood1944}, p.\,107, Titchmarsh~\cite{titchmarsh1939},

\noindent\textbf{§5.71.}\quad For Hardy's theorem, see Hardy~\cite{hardy1933}, and Titchmarsh~\cite{titchmarsh1937}, p.\,174, where further references are given.

\noindent\textbf{§7.}\quad Paley and Wiener~\cite{paleywiener1934} were the first to make a systematic study of Fourier transforms in the complex domain (one variable). The proof given here differs only in detail from the one presented by Dym and McKean~\cite{dymmckean1972}, Ch.\,3. For functions of exponential type see, for instance, Boas~\cite{boas1954}.

\noindent\textbf{§8.}\quad For generalities on Fourier orthogonal series see Kaczmarz and Steinhaus~\cite{kaczmarzsteinhaus1935}, Ch.\,II, where several examples of orthogonal systems are given, including Rademacher's~\cite{rademacher1922}, and Walsh's~\cite{walsh1923} which can be defined after Paley~\cite{paley1932}, in terms of Rademacher's functions. Let~$\chi_0(t)=1$, and if $N$ is a positive integer, expressed in the binary scale as $N=2^{n_1}+2^{n_2}+\cdots$, with $n_1>n_2>\cdots$, then $\chi_N(t)=\varphi_{n_1}(t)\cdot\varphi_{n_2}(t)$, where the~$(\varphi_n)$ are Rademacher's functions. The system~$(\chi_n)$ is orthonormal over~$(0,1)$, and complete, and is known as Walsh's.

\noindent Bessel's inequality (in several variables) has been shown by Siegel~\cite{siegel1935} to yield Minkowski's first theorem on lattice points in convex sets [cf.~the author's book~\cite{chandrasekharan1968}, p.\,99]. Atle Selberg~\cite{selbergtoappear} has shown that Bombieri's large sieve inequality can be viewed as a form of Bessel's inequality in a Hilbert space, cf.~H.L.~Montgomery~\cite{montgomery1978}.

\noindent For a concise introduction to Fourier trigonometric series in~$L_2(0,2\pi)$ see Hardy and Rogosinski~\cite{hardyrogo1944}. Series and integrals can be treated together on a group space. See Butzer and Nessel~\cite{butzernessel1971}. For Lebesgue's proof of the completeness of the trigonometric system, see Hardy and Rogosinski~\cite{hardyrogo1944}, or Zygmund~\cite{zygmund1959}, Vol.\,I, Ch.\,I, §6.

\noindent\textbf{§9.}\quad Hardy's~\cite{hardy1941} interpolation formula is also treated in Zygmund~\cite{zygmund1959}, Vol.\,II, p.\,276. As Dr.~Albert Stadler has remarked, the condition of boundedness in Theorem~\ref{thm:II-12} can be replaced by one of polynomial growth, in which case formula~\eqref{eq:II-9.7} will assume a more general form.

\noindent\textbf{§10.}\quad S.~Bernstein's work~\cite{bernstein1912} is also presented in Zygmund~\cite{zygmund1959}, Vol.\,II, p.\,11, Ch.\,X, and p.\,276, Ch.\,XVI. Zygmund's inequality for the integrated derivative of a trigonometric polynomial, as a generalization of Theorem~\ref{thm:II-14}, is given by him immediately after Bernstein's result. N.G.~de Bruijn~\cite{debruijn1947} has given generalizations of Bernstein's theorem for polynomials in the complex domain. For the Remark preceding Theorem~\ref{thm:II-14}, see Siegel~\cite{siegel1935}. See also Stein~\cite{stein1957}.

\section*{Notes on Chapter III}

\noindent\textbf{§11.}\quad For the extension of the Paley-Wiener theorem to $k$ dimensions, $k>1$, see Plancherel and P\'olya~\cite{plancherelpolya1937}, Stein~\cite{stein1957}, Stein and Weiss~\cite{steinweiss1971}, Ch.\,III, Th.\,4.9. The last-mentioned reference connects the theorem with the analysis of~$H^D$-spaces. See Narasimhan~\cite{narasimhan1968}, Ch.\,3, for the role of the Fourier transform in analytical problems on manifolds; also Ehrenpreis~\cite{ehrenpreis1970} in connexion with several complex variables.

\noindent\textbf{§1.}\quad For the basic theory of Stieltjes integrals see, for instance, Burkill and Burkill~\cite{burkillburkill1970}, Ch.\,6, and Widder~\cite{widder1941}, Ch.\,I.

\noindent\textbf{§2.}\quad As Zygmund has remarked, the essence of Theorems~\ref{thm:III-3} and~\ref{thm:III-4} is a classical result of the calculus of probability, in a form strengthened by Cram\'er. See Zygmund~\cite{zygmund1959}, Vol.\,II, Ch.\,XVI, Th.\,(4.24), p.\,262. See also Cram\'er~\cite{cramer1946}, Ch.\,10. Bochner has a generalization to~$E_k$, see Th.\,3.2.1 of his book~\cite{bochner1955}, p.\,56.

\noindent For Helly's theorem used in the proof of Theorem~\ref{thm:III-4}, see, for instance, Widder~\cite{widder1941}, Ch.\,I, §16, Th.\,16.2.

\noindent\textbf{§3.}\quad Theorems~\ref{thm:III-6} and~\ref{thm:III-7} are due to Bochner, see~\cite{bochner1932}, Th.\,23. He refers to previous work by F.~Bernstein and M.~Mathias. The generalization, without the assumption of continuity, is due to F.~Riesz~\cite{riesz1927}, who uses for the proof, however, his theorem on the representation of positive linear functionals, which is not used here in the proof of Theorem~\ref{thm:III-5}. For the Helly-Bray theorem used, see, for instance, Widder~\cite{widder1941}, p.\,31, Th.\,16.4. It is not necessarily true when the interval of integration is infinite, as Widder makes clear, hence the introduction of the kernel~$K_{R}(x)$. Carleman~\cite{carleman1944}, p.\,98, gives a proof of Bochner's theorem using the Poisson integral representation of functions which are positive and harmonic in a half-plane. A proof of the latter (see, for instance, Verblunsky~\cite{verblunsky1935}) can be obtained by using Herglotz's theorem~\cite{herglotz1911} on the representation of positive, harmonic functions in a circle (which is stated, for instance, in Stone~\cite{stone1932}, p.\,571), or more directly, as has been done by Loomis and Widder~\cite{loomiswidder1942} using the theorems of Helly, and of Helly-Bray. It should be remarked, however, that all these representation theorems are, more or less, of the same order of difficulty as Bochner's theorem, or Stone's spectral theorem~\cite{stone1932}, p.\,331, as was early recognized by F.~Riesz.

\noindent Apropos Corollary~\eqref{eq:III-3.27}, see Cram\'er~\cite{cramer1939}. For Theorem~\ref{thm:III-9} see Bochner~\cite{bochner1959}, p.\,329. Bochner has also a generalization to~$E_k$,~\cite{bochner1955}, Theorem~3.2.3, p.\,58.

\noindent For a generalization to distributions, see Schwartz~\cite{schwartz1950}, Vol.\,II, p.\,132, Th.\,XVIII; Schwartz makes a reference to Weil~\cite{weil1940}, p.\,122.

\noindent\textbf{§4.}\quad Lemma~\ref{lem:III-1} is due to H.A.~Schwarz. It is quoted by G.~Cantor, J.~f\"ur Math.~72 (1870), 141; and is given by Schwarz himself in his Ges.~Abhandlungen, II (1890), 341--343, with a reference to Cantor's quotation. Prof.~Raghavan Narasimhan has remarked that a rearrangement of Schwarz's argument is better adapted to generalization. ``If $f$ is real-valued, and continuous on $(a,b)$, and $\limsup_{\mathrm{h}\to0}(\Delta_{\mathrm{h}}^2f)(x)=0$, for all $x\in(a,b)$, where $\Delta_{\mathrm{h}}^2f(x)=h^{-2}\{f(x+h)+f(x-h)-2f(x)\}$, then $f$ is linear. To prove this, it is sufficient to prove that if $\limsup_{\mathrm{h}\to0}\Delta_{\mathrm{h}}^2f\geq0$, then $f$ is convex (i.e. if $\ell(x)=cx+d$, and $f(\alpha)\leq\ell(\alpha)$, $f(\beta)\leq\ell(\beta)$, where $a<\alpha<\beta<b$, then $f(x)\leq\ell(x)$ for $\alpha\leq x\leq\beta$, since one can argue similarly with $\liminf_{h\to0}\Delta_{\mathrm{h}}^2(-f)$. Since $\Delta_{\mathrm{h}}^2\ell=0$ for any linear function $\ell$, it is enough to prove that if $\limsup_{\mathrm{h}\to0}\Delta_{\mathrm{h}}^2f\geq0$, and $f(\alpha)\leq0$, $f(\beta)\leq0$, then $f(x)\leq0$ on $[a,\beta]$, i.e. that $f$ has no maximum on $(\alpha,\beta)$. Replacing $f(x)$ by $f(x)+\varepsilon x^2$, $\varepsilon>0$, one has only to show that if $\limsup_{\mathrm{h}\to0}\Delta_{\mathrm{h}}^2f>0$, on $(a,b)$, then $f$ has no local maximum on $(a,b)$. But this is obvious, for if $x_0$ is a local maximum, then $\Delta_{\mathrm{h}}^2f(x_0)\leq0$ for $h$ small enough, since $f(x_0)\geq f(x_0+h)$, $f(x_0)\geq f(x_0-h)$.'' Cf.~Narasimhan~\cite{narasimhan1985}, p.\,21--25.

\noindent Theorem~\ref{thm:III-10} is due to A.C.~Offord~\cite{offord1936}, and is the integral analogue of Cantor's fundamental theorem on the uniqueness of trigonometric series, which asserts that if a trigonometric series converges everywhere to zero, it vanishes identically; all its coefficients are zero. See, for instance, Zygmund~\cite{zygmund1959}, Vol.\,I, Ch.\,IX, p.\,326. Offord also proved~\cite{offord1937} a stronger theorem in which the hypothesis of convergence of the integral in~\eqref{eq:III-4.1} is replaced by $(C,1)$ summability. Offord shows that the stronger theorem is a ``best possible'', in the sense that even one exceptional point cannot be permitted, and $(C,1)$ summability cannot be relaxed to $(C,1+\varepsilon)$ summability for any $\varepsilon>0$. Zygmund's proof~\cite{zygmund1959}, Vol.\,II, Ch.\,XVI, §10, of Offord's first theorem is based on an equiconvergence theorem for trigonometric integrals and series which he treats in Vol.\,II, Ch.\,XVI, §9, and on results from Riemann's theory of trigonometric series which he treats in Vol.\,I, Ch.\,IX. For the use of equiconvergence theorems in analytic number theory, see, for instance, the author's book~\cite{chandrasekharan1970}, Ch.\,VIII. A generalization of Offord's stronger theorem to several variables would be of interest, though perhaps not easy. An equiconvergence theorem for trigonometric integrals in two variables has been given by H.~Keller~\cite{keller1976}. Functions of bounded variation in two variables come into play, and it is a moot question whether the notion of Vitali variation could be replaced by that of Fr\'echet, as Morse~\cite{morsetransue1950} and Transue did in another context.
